\section{Validation and reflection}
\label{sec:validation}

\subsection{The proposer and the judge}
\label{sec:validation-actor}

Acceptance has been the runtime's job since the chain, and nothing so far has asked whether the
runtime is any good at it. The three checks in \demoref{submit} are cheap and they are dull, and
both of those turn out to be load-bearing properties rather than shortcomings.

\begin{definition}[Actor and evaluator]
\label{def:actor-evaluator}
An agent separates \emph{proposal} from \emph{assessment} when the component that produces a candidate is not the component that scores it. Reflexion draws the split as two of three models: ``an Actor, denoted as $M_a$, which generates text and actions; an Evaluator model, represented by $M_e$, that scores the outputs produced by $M_a$'' \citep[\S3]{shinn2023reflexion}. The evaluator ``takes as input a generated trajectory and computes a reward score that reflects its performance within the given task context'', and it may be an exact-match grader, a heuristic written for the task, or another instance of a language model \citep[\S3]{shinn2023reflexion}.
\end{definition}

Self-Refine draws the same picture with one model in all three roles. It generates an output,
prompts the same model for feedback on that output, and prompts it again to revise, with three
prompts and no training \citep[\S2]{madaan2023selfrefine}. What matters for us is where its loop
exits, because the stopping condition ``either stops at a specified timestep $t$, or extracts a
stopping indicator (e.g.\ a scalar stop score) from the feedback''
\citep[\S2]{madaan2023selfrefine}, and the feedback is written by the model whose work is being
assessed.

Our loop exits on a value the runtime wrote. That difference does not show up in a diagram of boxes
and arrows, and it is not something a prompt can establish. \textbf{The accepting decision belongs
to whichever component writes the value the exit edge reads, which is a fact about the graph we
compiled and not about anything we asked the model to do.} Instruct a model to be strict with itself
and it remains the author of the number the edge tests.

So the question in front of us is not whether to separate the proposer from the judge, because
\demoref{submit} has been separate from the controller since the chain. The question is what a
second model adds to a separation we already have, and what it costs to let that model hold the
decision.

\subsection{What a model judge can and cannot see}
\label{sec:validation-judge}

Each check answers its question with a program rather than an opinion: i) the patch applies, so the
edit is well formed; ii) it touches the flagged line, so the run addressed its alert; iii) the
rescan is clean, so the scanner is quiet. Each is fast, each is reproducible,
and none of them needs a model.

Here, a verifier is \emph{sound} when it never accepts an incorrect solution, and \emph{complete}
when it never refuses a correct one. Soundness on its own buys nothing, because a verifier that
refuses everything never accepts an incorrect solution. \textbf{Completeness is what carries the
property the argument below rests on: a verifier that produces no false refusals leaves a system no
worse off than not checking at all, because anything generated correctly is still accepted.} The
external verifiers behind the gains in \cref{sec:validation-evidence} are both at once, being
decision procedures for their domains.

Our three checks are neither. A scanner that reports nothing has not established that nothing is
there, and an edit that renames the flagged expression or moves it behind a wrapper can clear an
alert while leaving the weakness in place, which is a false acceptance. A correct patch that fixes
the weakness somewhere other than the flagged line fails the second check, which is a false refusal.
What the checks are is cheap and reproducible, and the error to avoid is reading them as though they
were sound.

That gap is the validator's job, and it is a narrow one. Given the alert, the candidate patch, the
results of the three checks, and the trajectory that produced them, a validator can ask three
questions the checks cannot.

\bi

\item \textbf{Weakness or alert.} Does the patch remove the unsafe behavior, or does it remove the
scanner's ability to see it? This is the validation-retreat failure of \cref{sec:reactive-react}
under a different name, and on this fixture it is the failure most worth catching.

\item \textbf{Behavior.} Does the edit change what the endpoint returns? The performance measure
cares and none of the three checks looks.

\item \textbf{Provenance.} Did the run actually establish what its patch assumes, or did it assert
it? The trajectory is in front of the validator and was never in front of the checks.

\ei

\subsection{Why the judge does not decide acceptance}
\label{sec:validation-evidence}

Having found something for a second model to do, we now have to decide how much authority it gets,
and the evidence on that question is unusually one-sided.

Start with the setting the results are about. \emph{Intrinsic} self-correction is the case in which
a model revises ``based solely on its inherent capabilities, without the crutch of external
feedback'' \citep[\S1]{huang2024selfcorrect}, which is precisely the case where a validator holds
the accepting decision and nothing else does. Measured that way on reasoning benchmarks, accuracy
falls while cost rises: one model answers a grade-school mathematics benchmark at 95.5\% in a single
call, 91.5\% after one round of self-correction across three calls, and 89.0\% after two rounds
across five \citep[Table~3]{huang2024selfcorrect}. The same pattern appears on a commonsense
benchmark, where a weaker model falls from 75.8\% to 38.1\% after one round
\citep[Table~3]{huang2024selfcorrect}.

Supply the ground-truth label instead, so that correct answers are never revised, and the same
procedure improves that mathematics result from 95.5\% to 97.5\%
\citep[Table~2]{huang2024selfcorrect}. The gain is real and it is not available: ``If we are already
in possession of the ground truth, there seems to be little reason to deploy LLMs for
problem-solving'' \citep[\S3]{huang2024selfcorrect}. Two systems whose reported improvements rest on
that setup are named, and one of them is Reflexion \citep[Table~1]{huang2024selfcorrect}, which
matters for how we read \cref{sec:reflection}.

A second study separates the verifier from the generator and varies only the verifier. Across 100
instances per domain with one model, self-critique moves a graph-coloring result from 16\% to 2\%
and a Game of 24 result from 5\% to 3\%. Replacing the model verifier with a sound one moves the
same two results to 38\% and 36\% \citep[Table~1]{stechly2025selfverification}. On a
blocks-planning domain self-critique does help, from 40\% to 55\%, and a sound verifier still beats
it at 60\% to 87\% depending on how much detail the critique carries
\citep[Table~1]{stechly2025selfverification}.

The verifier's error profile explains the collapse. Used as a verifier, the model reaches 72.4\%
accuracy on graph coloring, refusing 95.8\% of the correct answers it saw and accepting 6.5\% of the
incorrect ones, and 79.6\% on a disguised blocks domain, where those two rates are 97.09\% and 0.5\%
\citep[Table~2]{stechly2025selfverification}. A judge
that refuses nearly every correct answer does not merely fail to help, because ``the system rejects
most answers and then times out on a set of later, worse generations''
\citep[\S5]{stechly2025selfverification}. The acceptance error runs the other way. On the two
domains where the judge does least badly it accepts 10.4\% of incorrect Game of 24 answers and
18.55\% of incorrect blocks plans, against 6.5\% and 0.5\% on the two where it refuses almost
everything \citep[Table~2]{stechly2025selfverification}.

One finding sets our design rule. Once the verifier is sound, binary feedback, feedback naming the
first error, and feedback naming every error land within four points of one another on three of the
four domains. Those readings are 38\%, 37\% and 34\% on graph coloring, 36\% against 38\% on Game of
24, and 10\%, 8\% and 6\% on the disguised blocks domain. On blocks planning they do not, running 60\%, 87\% and
83\%, so there naming the first error is worth twenty-seven points over a bare refusal
\citep[Table~1]{stechly2025selfverification}. \textbf{What a loop needs first from its judge is a
trustworthy accept-or-reject bit; what an attached explanation adds on top of that ran from nothing
to twenty-seven points across the four domains.} Dropping the critique altogether and sampling
repeatedly does at least as well on two of them, reaching 44\% on graph coloring and 42\% on Game of
24 at twenty-five samples \citep[Table~1]{stechly2025selfverification}.

These are four reasoning and planning domains with one model family, not program repair, and repair
has something those domains lack: a rescan and a test suite are real programs that run. What
transfers is the error profile of the source rather than a rate we can borrow. Where the accepting
signal comes from a model, the outcome follows that model's error rates rather than the idea. The
loop collapses on the two domains whose false-refusal rates run above 95\%, and on blocks planning, where that rate is
15.48\%, it gains fifteen points where a decision procedure gains twenty to forty-seven
\citep[Tables~1 and 2]{stechly2025selfverification}. Where the signal comes from a verifier that is
correct in both directions, the loop improves on all four. Our three checks are programs and they
are correct in neither direction, so what they inherit is reproducibility rather than the
guarantee.

That is why the mechanical checks keep the accepting decision. \textbf{The validator may refuse a candidate that
passed all three checks, and it may not release one that failed them.} Confined that way, the only
error it can make is refusing work that was already finished, and a refusal costs a turn against a
budget we set. Were it allowed to release, its false acceptances would end the run, and nothing
downstream would look again. Given false-negative rates of the size measured above, a judge confined
to refusal can still stall a run by refusing everything, which is what the attempt cap of
\cref{sec:reflection} is for.

\Cref{sec:fixed} and \cref{sec:reactive} both promised that acceptance would move to something that
did not write the patch, and it has. \demoref{submit} is that something, and the validator is a
second one standing beside it. Neither section promised that the judge would be a model, and the
evidence above is why we do not make it one.

\subsection{What survives a refusal}
\label{sec:reflection}

A refusal ends an attempt, and it need not end everything the attempt learned. What crosses that
boundary is a design decision, and the two obvious answers are both wrong.

The loop itself is short. The validator refuses a candidate, the attempt limit permits another try,
and the runtime starts the next attempt from a cleared record. Something has to change between the
two, or the second attempt is the first one repeated at twice the price. What changes is the
context the next attempt begins with, and we get to choose what goes into it.

\bi

\item \textbf{The whole trajectory.} Too large to carry, and it is the accumulation we were trying
to get out of. An attempt that inherits every observation of the attempt before it is a longer first
attempt wearing a new label.

\item \textbf{The patch.} The artifact that just failed. Carrying it invites the next attempt to
defend it, which is the one thing we do not want a fresh attempt doing.

\item \textbf{The thought.} The action that changes the context and touches nothing else
(\cref{sec:reactive-loop}). It costs a model call, it has no effect in the environment, and it is
the only one of the three that is smaller than what produced it.

\ei

\begin{definition}[Reflection step]
\label{def:reflection}
A reflection step converts the outcome of a failed attempt into text that conditions later attempts.
Reflexion adds it as the third of its three models, ``a Self-Reflection model, denoted as $M_{sr}$,
which generates verbal reinforcement cues to assist the Actor in
self-improvement'' \citep[\S3]{shinn2023reflexion}. Given ``a sparse reward signal, such as a binary
success status (success/fail), the current trajectory, and its persistent memory'', it produces
feedback that ``is then stored in the agent's memory'' \citep[\S3]{shinn2023reflexion}.
\end{definition}

Two lifetimes come with that arrangement, and they are the same two our state record keeps: ``the
trajectory history serves as the short-term memory while outputs from the Self-Reflection model are
stored in long-term memory'' \citep[\S3]{shinn2023reflexion}. The trajectory dies with the attempt;
the reflection does not.

Our reflection node runs on rejection and writes one short entry from four inputs: the alert, the
candidate patch, the validator's reasons, and the trajectory that produced them. The entry appends
to \demoref{reflections}, which survives the reset, while the observations, the plan, the candidate
and the validation are cleared for the new attempt. The counters survive too, so the budget of forty
model calls still binds across both attempts rather than restarting with each one.
\textbf{The reset is what makes the guidance worth anything, because guidance delivered alongside
the context that produced the failure is just more context.} Which fields are cleared, and why
clearing an appending field takes an explicit assignment, is the state machinery of Part~B.

What this does not buy is a second opinion from nowhere. The signal our reflection step receives is
the validator's rejection reasons, which are the three mechanical checks plus a model's objections,
and \cref{sec:validation-judge} said what those checks are worth: they are neither sound nor
complete. The measurements in \cref{sec:validation-evidence} were about revision with no external
signal at all, and this is not that case, but it is not the ground-truth case either. What we expect
is improvement where the failure was legible in the reasons the run was given, and no improvement at
all on an alert for which no patch exists. The attempt cap is there because neither outcome
announces itself.

\subsection{What we gain and what we lose}
\label{sec:validation-tradeoff}

A validator and a reflection step are additions of the same kind as the planner: we wrote the
validator's authority, the routing predicate and the reset contract, and every one of those is
still design time. What changes is that a second model now speaks about work the first model did.

\subsubsection{What we gain}

\bi

\item \textbf{Refusal.} A patch that passes all three mechanical checks can still be turned back,
which closes the one gap those checks were always going to leave open.

\item \textbf{Memory.} A refused attempt leaves guidance behind, so a second attempt starts from
what the first one was told rather than from the empty context the first one started in.

\ei

\subsubsection{What we lose}

\bi

\item \textbf{Judgement.} A second model is a second source of error, and its refusals are
expensive in a way its silence is not. A validator that refuses good work spends the budget on work
that was already done.

\item \textbf{Attempts.} The attempt loop multiplies everything beneath it, because a second attempt
pays for a second plan, a second set of turns and a second validation against the same budget. The
guidance it starts from was written by the model whose attempt failed.

\ei

What \cref{sec:planning-tradeoff} left open was that nothing checked the finished work, and a
validator closes it without returning any of the guarantees \cref{sec:fixed-tradeoff} named. Shape
and termination stay gone, the attempt cap replaces termination with a counter we set, and the
accepting decision stays with the three checks, because neither the validator nor the reflection
step is safe to hold it.

\subsection{The demo: the validator and the reflection step}
\label{sec:validation-demo}

\begin{demostub}
Expected artifact: the validation node and the reflection node of
\demoref{demo/lec2\_planning\_and\_validation\_agent}, with one recorded run of each under
\demoref{demo/runs/}, beside the planner artifacts of \cref{sec:planning-demo}.

For the validator, one selected candidate that passes all three checks in \demoref{submit} and is
refused anyway, with the four inputs it was given and the reasons it returned. A run in which the
mechanical checks fail and the validator is never consulted belongs beside it, because the order of
the two is the design rule of \cref{sec:validation-evidence}.

For the reflection step, one rejected attempt and the attempt that follows it, side by side. Four
things go on the page: the reasons the validator returned, the entry appended to
\demoref{reflections}, the fields the runtime cleared, and the first model input of the second
attempt, where the guidance arrives without the context that produced it. The usage ledger runs
across both attempts in the same view.
\end{demostub}

\subsection{Exercises}

\be

\item Give the validator authority to release a patch as well as refuse one. Using a false-positive
rate of the order measured in \cref{sec:validation-evidence}, say what the worst case costs on a
batch of forty alerts, and why the mechanical checks are the cheaper place to keep the decision.

\item A rejected attempt leaves one entry in \demoref{reflections}. Write the entry you would want
after a patch that cleared the scanner alert by renaming the flagged expression, and say which of
the validator's three questions your entry would have to have been told about.

\item Argue the other side of the design rule in \cref{sec:validation-evidence}: describe a fixture
on which letting the validator release a patch is the right choice, and name the property of that
fixture that ours does not have.

\ee

\subsection{Reading}

\citet{madaan2023selfrefine}, \S2 and Algorithm~1, for a refinement loop in which one
model fills every role. Then \S3 of \citet{shinn2023reflexion} for all three of its models, the
actor and evaluator split of \cref{sec:validation-actor} and the self-reflection model and memory of
\cref{sec:reflection}.

The counterweight is the part of the reading to do slowly. \citet{huang2024selfcorrect}, \S\S1 and 3
with Tables~1 to 3, defines intrinsic self-correction and measures it. \citet{stechly2025selfverification},
\S5 with Tables~1 and 2, separates the verifier from the generator and varies only the verifier.
Read both against the claim \cref{sec:validation-actor} makes about who writes the value the exit
edge reads.

\paragraph{Looking ahead.}
The serial agent is now complete. A plan gives the run something to fall short of, the mechanical
checks and the validator give it a refusal it cannot argue with, and reflection gives a refused
attempt something to leave behind. Every version of it has been one model call at a time, on one
trajectory, against one alert, and each addition was paid for in a guarantee we gave up.
\Cref{sec:patterns} asks what changes when the work is divided among several agents instead, and
what each division costs in the same currency.
